\section{Discussion and Conclusion}

This literature review has examined the evolution of power delivery architectures for sub-5\,nm CMOS, from conventional front-side PDNs through buried power rails to full backside delivery. The reviewed studies consistently demonstrate that BSPDNs provide substantial improvements in power integrity and design efficiency. By relocating the power network to the wafer backside, this architecture decouples power and signal routing, enabling IR~drop reductions of more than four times\cite{8932458}, voltage droop improvements of 30\%\cite{10185208}, frequency gains of 6\%\cite{10185208}, and cell area savings of up to 19\%\cite{10631463}. These results, drawn from both simulation and fabricated test vehicles, indicate that backside delivery addresses the fundamental scaling limitations that front-side PDNs face at advanced nodes. Buried power rails serve as an effective intermediate step, reducing IR~drop by 25\% and voltage droop by 17\%\cite{9817394}, but the full separation of power and signal routing achieved by BSPDNs delivers substantially greater benefits across all evaluated metrics.

However, these gains come with trade-offs that cannot be overlooked. The thinned substrate required for nano-TSV connections removes lateral heat-spreading capability, leading to hotspot temperatures up to 45\% higher than equivalent front-side designs\cite{10925422}. Self-heating effects further complicate the picture at the 3\,nm node, where power-thermal co-design becomes essential to avoid reliability degradation\cite{10529350}. Additionally, the reduced DC resistance of BSPDNs shifts impedance resonance peaks into frequency ranges that may coincide with circuit switching activity\cite{8993617}, requiring careful decoupling capacitor placement during design. These challenges suggest that the electrical benefits of BSPDNs cannot be evaluated in isolation from their thermal and dynamic consequences.

This review has several limitations that should be acknowledged. The analysis is based on nine sources drawn exclusively from IEEE Xplore, which may not capture relevant work published in other databases or in non-English venues. Furthermore, the majority of the reviewed results are based on simulation rather than measured silicon data; only Intel's PowerVia study\cite{10185208} reports results from a fabricated test vehicle. The scope of the review is also limited to power integrity and routing metrics, and does not address manufacturing yield, cost, or supply chain considerations that will influence industrial adoption.

Future research should focus on three areas. First, measured thermal and electrical data from production-grade BSPDNs are needed to validate the simulation-based predictions that dominate the current literature. Second, power-thermal co-design methodologies that jointly optimise IR~drop and hotspot temperature at the block and system level remain an open challenge. Third, the extension of backside routing to include signal interconnects\cite{11192533} opens new possibilities for standard-cell scaling, but the associated design rules and parasitic trade-offs require further investigation. Addressing these questions will be critical to realising the full potential of backside power delivery in future technology nodes.
